\documentclass[a4paper]{article}
% Español (México) con XeLaTeX: fontspec para fuentes Unicode, polyglossia
% para la división silábica y los nombres del documento en la variante mexicana.
\usepackage{fontspec}
\usepackage{polyglossia}
\setdefaultlanguage[variant=mexican]{spanish}
% Los nombres chinos van como «pinyin (original)»: xeCJK da a los caracteres
% Han su propia fuente; sin ella XeLaTeX los omite con un aviso.
% FandolSong no cubre algunos caracteres de nombres (U+73FA); WenQuanYi Zen Hei sí.
\usepackage[AutoFallBack=true]{xeCJK}
\setCJKmainfont{FandolSong-Regular.otf}
\setCJKfallbackfamilyfont{\CJKrmdefault}{WenQuanYi Zen Hei}
% polyglossia no traduce los nombres de listings.
\AtBeginDocument{\providecommand{\lstlistingname}{}\renewcommand{\lstlistingname}{Listado}%
  \providecommand{\lstlistlistingname}{}\renewcommand{\lstlistlistingname}{Índice de listados}}
\usepackage{amsmath,amssymb}
\usepackage{graphicx}
\usepackage[margin=2.35cm]{geometry}
\usepackage[most]{tcolorbox}
\usepackage{etoolbox}
\usepackage{listings}
\usepackage{xcolor}
\usepackage{hyperref}
\usepackage{booktabs,longtable,tabularx,array}
\usepackage{subcaption}
\usepackage{float}
\usepackage{enumitem}
\usepackage{microtype}
\usepackage{fancyhdr}
\usepackage{needspace}

\definecolor{kbBack}{HTML}{F2F6FA}
\definecolor{kbFrame}{HTML}{9DB4CC}
\definecolor{kbTitle}{HTML}{52708F}
\definecolor{ibBack}{HTML}{FBF5E7}
\definecolor{ibFrame}{HTML}{C9A961}
\definecolor{ibTitle}{HTML}{7D5F1E}
\definecolor{wbBack}{HTML}{FCF2F0}
\definecolor{wbFrame}{HTML}{D6A096}
\definecolor{wbTitle}{HTML}{9E4F43}
\definecolor{dbBack}{HTML}{F4F7F3}
\definecolor{dbFrame}{HTML}{AEC2AC}
\definecolor{dbTitle}{HTML}{5F7F64}

\tcbset{softbox/.style={
  enhanced,breakable,boxrule=0.8pt,arc=2pt,left=3mm,right=3mm,
  top=2mm,bottom=2mm,coltitle=white,fonttitle=\bfseries,
  toptitle=0.6mm,bottomtitle=0.6mm
}}
\newtcolorbox{knowledgebox}[1]{softbox,colback=kbBack,colframe=kbFrame,colbacktitle=kbTitle,title={#1}}
\newtcolorbox{importantbox}[1]{softbox,colback=ibBack,colframe=ibFrame,colbacktitle=ibTitle,title={#1}}
\newtcolorbox{warningbox}[1]{softbox,colback=wbBack,colframe=wbFrame,colbacktitle=wbTitle,title={#1}}
\newtcolorbox{dialoguebox}[1]{softbox,colback=dbBack,colframe=dbFrame,colbacktitle=dbTitle,title={#1}}

\lstset{
  language=Python,basicstyle=\ttfamily\small,
  keywordstyle=\color{kbTitle}\bfseries,stringstyle=\color{wbTitle},
  commentstyle=\color{dbTitle},breaklines=true,frame=single,numbers=left,
  numberstyle=\tiny\color{gray},captionpos=b,columns=fullflexible,
  keepspaces=true,showstringspaces=false,extendedchars=true
}
\BeforeBeginEnvironment{lstlisting}{\Needspace{12\baselineskip}}

\hypersetup{colorlinks=true,linkcolor=kbTitle,urlcolor=kbTitle,citecolor=wbTitle}
\setlist{itemsep=0.25em,topsep=0.4em}
\setlength{\parindent}{2em}
\setlength{\parskip}{0.25em}
\newcolumntype{L}[1]{>{\raggedright\arraybackslash}p{#1}}

\providecommand{\notetitle}{Notas del curso: [tema del curso]}
\providecommand{\noteauthors}{Elaboradas a partir de los materiales públicos de Stanford CS25}
\providecommand{\notedate}{Versión reescrita en 2026}
\providecommand{\videochannel}{Stanford Online}
\providecommand{\videopublishdate}{2022}
\providecommand{\videoduration}{[duración del video]}
\providecommand{\videourl}{}
\providecommand{\videocoverpath}{cover.jpg}
\providecommand{\coursesite}{https://web.stanford.edu/class/cs25/}
\providecommand{\repourl}{https://github.com/hqhq1025/ai-course-notes}
\providecommand{\skillversion}{wdkns/wdkns-skills@39f1a04}

\newcommand{\figsource}[1]{\par\vspace{-0.3em}{\footnotesize\color{gray}\emph{Fuente: #1}}\par}
\newcommand{\teachervoice}[1]{\begin{knowledgebox}{Nota de clase}#1\end{knowledgebox}}
\newcommand{\term}[1]{\textbf{#1}}

\pagestyle{fancy}
\fancyhf{}
\fancyfoot[C]{\thepage}
\fancyfoot[R]{\footnotesize\color{gray}\href{\repourl}{AI Course Notes}}
\renewcommand{\headrulewidth}{0pt}
\renewcommand{\footrulewidth}{0pt}

\newcommand{\makecscover}{
\begin{titlepage}
\centering
\vspace*{0.7cm}
{\Large Stanford CS25 · Transformers United\par}
\vspace{0.8cm}
{\huge\bfseries \notetitle\par}
\vspace{0.6cm}
{\large \noteauthors\par}
\vspace{0.25cm}
{\large \notedate\par}
\vspace{0.9cm}
\ifdefempty{\videocoverpath}{}{
  \includegraphics[width=0.86\textwidth,height=0.43\textheight,keepaspectratio]{\videocoverpath}\par
}
\vfill
\begin{tcolorbox}[width=0.92\textwidth,colback=black!2!white,colframe=black!45,boxrule=0.7pt,arc=2pt]
\textbf{Autor o canal del video}: \videochannel\par
\textbf{Fecha de publicación}: \videopublishdate\par
\textbf{Duración del video}: \videoduration\par
\textbf{Sitio del curso}: \href{\coursesite}{\nolinkurl{\coursesite}}\par
\ifdefempty{\videourl}{}{\textbf{Enlace del video}: \href{\videourl}{\nolinkurl{\videourl}}\par}
\textbf{Normas de elaboración}: \skillversion\ + las normas source-first / teacher-voice / visual-QA de este repositorio
\end{tcolorbox}
\end{titlepage}
\tableofcontents
\newpage
}

\usepackage{multicol}

\renewcommand{\notetitle}{Lecture 23: No Language Left Behind}
\renewcommand{\noteauthors}{Elaborado a partir de la clase de Angela Fan, 39 estados de slide de la clase, subtítulos hechos por personas y materiales oficiales de NLLB}
\renewcommand{\notedate}{CS25 V3 · versión reescrita source-first de 2026}
\renewcommand{\videochannel}{Stanford Online}
\renewcommand{\videopublishdate}{Clase: 2023-11-14; publicación: 2023-12-15}
\renewcommand{\videoduration}{52:28}
\renewcommand{\videourl}{https://www.youtube.com/watch?v=ckNMsUuLryM}
\renewcommand{\videocoverpath}{cover.jpg}

\newcommand{\lecturefigure}[3]{%
\begin{figure}[H]
  \centering
  \includegraphics[width=0.96\textwidth,height=0.54\textheight,keepaspectratio]{slides-images/#1}
  \caption{#2}
  \figsource{#3}
\end{figure}
}

\begin{document}
\makecscover

\section{Introducción: doscientos idiomas no son una cifra de un modelo}

Esta clase no trata de «escalar un poco más el Transformer», sino de un problema de sistemas más difícil: cuando los datos disponibles, las tradiciones de escritura, los estándares lingüísticos, los recursos de evaluación y las necesidades reales están muy desequilibrados, ¿cómo se extiende la traducción automática de unos pocos idiomas de muchos recursos a doscientos idiomas, sin perder calidad, seguridad ni utilidad? NLLB (No Language Left Behind) divide este objetivo en una cadena completa: primero entrevista a los usuarios, luego construye conjuntos de evaluación y datos semilla, después extrae corpus mediante minería, entrena modelos multilingües y contiene el sobreajuste, y por último valida el resultado con métricas automáticas, evaluación humana y revisiones de toxicity.

\lecturefigure{slide-01-title.jpg}{No Language Left Behind: de la cobertura de idiomas a una traducción automática centrada en las personas}{Video de la clase de Stanford 00:00:54}

\begin{knowledgebox}{Lectura de la figura: las tres palabras clave del título}
\term{No Language Left Behind} no garantiza que todos los idiomas ya tengan la misma calidad; es una línea de investigación. \term{Scaling} no se refiere solo a la cantidad de parámetros, sino también al número de idiomas, a las direcciones de traducción, al pipeline de datos y a la capacidad de evaluación. \term{Human-Centered} exige que usuarios reales participen en la definición de necesidades, los estándares de datos, la validación humana y los límites de seguridad. Toda la clase se organiza en torno a estos tres niveles.
\end{knowledgebox}

\teachervoice{Al iniciar, Fan comentó que también podría haber hablado de Llama 2, pero que NLLB está igual de relacionado con el tema de IA generativa del curso y que, además, este problema tiene un significado personal para ella: el inglés es su tercer idioma. Esta motivación nos recuerda que los sistemas multilingües no son un benchmark abstracto, sino que determinan a quién puede servir la tecnología.}

\begin{importantbox}{La tesis de sistemas de esta clase}
La calidad de la traducción automática multilingüe no es función de un solo modelo, sino de un sistema acoplado:
\[
Q = F(H,E,S,M,V),
\]
donde $H$ representa las human needs y los mecanismos de participación, $E$ la evaluation infrastructure, $S$ los source data y el mining pipeline, $M$ el model/training y $V$ la validation y la safety. Si falta cualquiera de estos eslabones, aumentar los parámetros o el volumen de datos puede limitarse a amplificar los sesgos existentes.
\end{importantbox}

\subsection{Ruta de lectura}

Esta clase explica primero por qué «cuántos idiomas se admiten» no basta por sí solo para medir el avance; después, siguiendo el orden de la clase, analiza las entrevistas, FLORES-200, NLLB-Seed, WikiMatrix/LASER3/Stopes, la receta de datos, MoE y el curriculum learning; por último, integra las métricas automáticas, el juicio humano, la toxicity y las preguntas y respuestas de la clase (Q\&A) en un marco de evaluación que admite auditoría.

\subsection{Resumen de la sección}

El verdadero objeto de investigación de NLLB es el «sistema de cobertura de idiomas», no un modelo aislado. Cada conjunto de datos, arquitectura y métrica que aparece más adelante debe entenderse dentro del ciclo completo de necesidades, evidencia y seguridad.
\section{De la cobertura lingüística a la utilidad real}

Para entender el punto de partida de NLLB, primero hay que distinguir tres cosas que a menudo se confunden: cuántas lenguas escritas existen en el mundo, cuántas lenguas cubrían los productos comerciales en ese momento y cuántas de ellas alcanzaban realmente una calidad de traducción utilizable, confiable y segura. En clase, una serie de slides muy sencillas desarrolla estos tres niveles uno por uno.

\subsection{Tres mil lenguas escritas y una motivación personal}

El primer mapa del mundo no busca un número total de lenguas exacto e invariable, sino establecer un orden de magnitud: las lenguas escritas del mundo superan con mucho lo que cubren los principales productos digitales. También sugiere un problema más profundo: «tener escritura» no equivale a «tener abundante texto digital», y menos aún a contar con corpus paralelos de alta calidad que puedan usarse directamente para entrenar.

\lecturefigure{slide-02-written-languages.jpg}{En el mundo existen más de tres mil lenguas escritas}{Video de la clase de Stanford 00:01:24}

\begin{knowledgebox}{Lectura de la figura: existencia de la lengua, digitalización y aptitud para el entrenamiento son tres condiciones distintas}
Las lenguas del mapa necesitan, en primer lugar, un identificador estable de la lengua o de la variante; en segundo lugar, texto digital accesible; y solo en tercer lugar pueden llegar a tener corpus bilingües alineados a nivel de oración. Una lengua de bajos recursos puede estar muy viva en el primer nivel y ser casi invisible en los otros dos. Por lo tanto, tener pocos datos de entrenamiento no es sinónimo de tener pocos hablantes, de que la lengua valga menos ni de que la lengua sea simple.
\end{knowledgebox}

\begin{warningbox}{No tomar el «número de lenguas escritas» como un valor verdadero fijo}
Los límites entre lengua y dialecto, los sistemas de escritura y las formas de normalización implican decisiones sociales. La cifra de la clase sirve para mostrar la brecha de cobertura; no debe usarse para asignar a la fuerza un código único a cada lengua ni para negar el criterio de las comunidades locales sobre nombres, escrituras y variantes.
\end{warningbox}

\teachervoice{El docente enfatiza que la tecnología multilingüe no es algo que nunca haya existido; la traducción automática es, de hecho, una de las aplicaciones comerciales de la IA generativa más exitosas y extendidas. La cuestión no es «si hay traducción», sino por qué esa capacidad sigue concentrada en unas pocas lenguas.}

\subsection{La cobertura de un producto es una instantánea en el tiempo, no una prueba de calidad}

A continuación, la clase coloca en una misma figura el número de lenguas del mundo y la cobertura que tenían en ese momento Google Translate y Microsoft Translator. El primer propósito de esta comparación es mostrar la diferencia en cantidad; el segundo, recordar que el conteo de lenguas en la página de un producto solo indica que la interfaz las admite, y no dice nada sobre la distribución de errores en textos largos, dominios especializados, direcciones de bajos recursos o escenarios donde la seguridad importa.

\lecturefigure{slide-03-product-coverage.jpg}{Cobertura lingüística de los productos comerciales de traducción en la fecha de la clase}{Video de la clase de Stanford 00:02:03}

\begin{knowledgebox}{Lectura de la figura: primero comparar órdenes de magnitud, después preguntar cómo se define la calidad}
En la figura, la cobertura de los productos es del orden de cien lenguas, mientras que las lenguas escritas son del orden de miles. A primera vista debe notarse la diferencia de cantidad; en una segunda lectura hay que preguntar si cada «lengua admitida» incluye traducción en ambos sentidos, si fue validada por hablantes nativos y si funciona igual de bien en las direcciones de bajos recursos. El número de lenguas es un indicador de coverage, no de quality, safety ni equity.
\end{knowledgebox}

\teachervoice{Al dar el número de lenguas de los productos, Fan aclaró por iniciativa propia que esas cifras «quizá ya estén un poco desactualizadas». Por eso, estas notas conservan como referencia temporal la clase de 2023 y no presentan las cifras de la slide como especificaciones vigentes de los productos en 2026.}

\subsection{El objetivo no es duplicar, sino duplicar de forma segura}

Esta sección convierte la diferencia de cantidad anterior en un objetivo de ingeniería. La pregunta del proyecto NLLB se resume en una frase: si se parte de unas cien lenguas, ¿cómo duplicar la cobertura? Pero el paréntesis que la slide agrega después es el núcleo de la restricción: no basta con añadir etiquetas de lengua; también hay que obtener, en la medida de lo posible, traducciones de alta calidad, seguras y realmente utilizables.

\lecturefigure{slide-04-double-coverage-safely.jpg}{Duplicar la cobertura lingüística manteniendo alta calidad y seguridad}{Video de la clase de Stanford 00:02:51}

\begin{importantbox}{Cómo formular correctamente el objetivo de cobertura}
Sea $\mathcal{L}$ el conjunto de lenguas y $\mathcal{D}\subseteq\mathcal{L}\times\mathcal{L}$ el conjunto de direcciones que el sistema puede afirmar que cubre. El objetivo real no es maximizar solamente $|\mathcal{L}|$ o $|\mathcal{D}|$, sino cumplir en cada dirección restricciones mínimas de calidad, seguridad y utilidad:
\[
\max |\mathcal{D}|\quad
\text{s.t.}\quad q_d\ge \tau_q,\ s_d\ge \tau_s,\ u_d\ge \tau_u\quad \forall d\in\mathcal{D}.
\]
Aquí $q_d$ representa la calidad de la traducción, $s_d$ el desempeño en seguridad y $u_d$ la utilidad para los usuarios; $\tau$ es el umbral mínimo que define el proyecto. La clase no dio umbrales concretos; la fórmula sirve para ilustrar las restricciones de múltiples objetivos, no para inventar reglas de gobernanza.
\end{importantbox}

\begin{warningbox}{La cobertura lingüística no se cuenta por etiquetas}
Una misma lengua puede abarcar distintas escrituras, variantes regionales y dominios; un mismo par de lenguas se divide además en direcciones into-English, out-of-English y non-English directions. Si solo se agregan nombres de lenguas a un menú desplegable, el sistema puede seguir sin ser utilizable justo en las direcciones donde más se necesita.
\end{warningbox}

\subsection{«Altos recursos» describe los datos, no la población}

Al repasar la evolución histórica, la clase señala que la traducción automática se desarrolló durante mucho tiempo en torno a las lenguas de altos recursos, y que corpus paralelos institucionales como Europarl fueron la base inicial. El término high-resource describe principalmente los datos de entrenamiento y evaluación disponibles, no la población que habla la lengua. Una lengua con muchísimos hablantes puede seguir siendo de bajos recursos por falta de texto digital, de normalización o de corpus paralelos.

\lecturefigure{slide-05-current-progress.jpg}{La traducción automática pasa gradualmente de las lenguas de altos recursos a las de bajos recursos}{Video de la clase de Stanford 00:04:00}

\begin{knowledgebox}{Lectura de la figura: el eje de recursos y el eje de población son independientes}
La trayectoria histórica de la izquierda depende de textos bilingües institucionales; los avances recientes de la derecha dependen del modelado multilingüe, la supervisión débil y la minería de corpus. Al leer la figura conviene separar el «volumen de recursos de datos» del «tamaño de la población»: tener muchos recursos suele facilitar el modelado, pero no indica una mayor necesidad social; tener pocos recursos debilita a la vez el entrenamiento, el ajuste de hiperparámetros y la evaluación.
\end{knowledgebox}

\begin{knowledgebox}{Términos clave: recursos y datos}
\begin{tabularx}{\textwidth}{@{}L{0.21\textwidth}L{0.28\textwidth}X@{}}
\toprule
Término & Problema que resuelve & Significado en esta clase \\
\midrule
High-resource language & Describir la suficiencia de datos & Cuenta con abundantes recursos monolingües, bilingües y de evaluación; no equivale a tener más hablantes \\
Low-resource language & Describir la escasez de datos & Faltan recursos de entrenamiento, de normalización y de evaluación humana \\
Monolingual data & Aprender la distribución de una sola lengua & Textos, como páginas web o documentos, que contienen una sola lengua \\
Bitext / parallel data & Aprender la correspondencia entre lenguas & Pares de oraciones de origen y destino alineadas semánticamente \\
Translation direction & Precisar el sentido de entrada y salida & $A\rightarrow B$ y $B\rightarrow A$ son dos tareas distintas \\
\bottomrule
\end{tabularx}
\end{knowledgebox}

\subsection{Resumen de la sección}

El punto de partida de NLLB no es «en el mundo hay tres mil lenguas, así que hay que entrenar tres mil», sino incorporar en el objetivo, de forma conjunta, la cobertura, la calidad, la seguridad y la utilidad. Los altos y bajos recursos describen las condiciones de los datos; el número de lenguas de un producto es solo una instantánea en el tiempo y no sustituye la evaluación por dirección.
\section{Primero preguntar a las personas, después definir la tarea}

Si el objeto de estudio está directamente ligado a la vida comunitaria, la educación y la continuidad de la lengua, el equipo del modelo no puede imaginar las necesidades basándose solo en un benchmark. NLLB primero confirmó el problema mediante entrevistas y luego hizo participar a hablantes nativos en la traducción y la evaluación; este es el giro metodológico más importante de la clase, y también el que con más facilidad queda oculto tras los detalles del modelo.

\subsection{Cómo saber que se está resolviendo un problema real}

Esta sección pasa de los objetivos de cobertura a la definición del problema. Antes de mostrar cualquier pipeline de datos, Fan plantea un criterio de investigación: lo más importante es confirmar que el equipo está resolviendo un problema real, sobre todo cuando la tecnología está muy cerca de las personas. Esta pregunta no es un lema ético, sino el primer paso que determina los estándares de datos, las condiciones de éxito y el costo de los errores en lo que sigue.

\lecturefigure{slide-06-real-problem.jpg}{Confirmar, antes de modelar, que el problema existe para los usuarios reales}{Video de la clase de Stanford 00:05:06}

\begin{knowledgebox}{Lectura de la figura: es una compuerta del proceso de investigación, no una figura retórica de apertura}
Si a los usuarios les preocupa sobre todo la pérdida de la lengua en la educación, el sistema no puede optimizar solo la traducción de noticias; si la comunidad usa varios sistemas de escritura o variantes locales, unificarlos en un único estándar puede destruir directamente la utilidad. Por eso, las entrevistas sobre necesidades modifican el muestreo de datos, el sistema de etiquetas, los escenarios de evaluación y la release policy.
\end{knowledgebox}

\begin{importantbox}{Human-centered no equivale a hacer una prueba con usuarios al final}
La cadena de participación más completa es:
\[
\text{Entrevistas sobre necesidades}\rightarrow\text{Negociación del estándar lingüístico}\rightarrow\text{Traducción profesional}
\rightarrow\text{Evaluación por hablantes nativos}\rightarrow\text{Retroalimentación y corrección de errores}.
\]
Los usuarios no son inspectores de aceptación al final del pipeline, sino copartícipes en la definición del problema y en la producción de evidencia.
\end{importantbox}

\subsection{Muestra de entrevistas y límites del método}

El equipo entrevistó a 44 hablantes nativos de 36 lenguas y de varias regiones. Esta muestra no puede representar a todas las comunidades de lenguas de bajos recursos, pero basta para exponer problemas que una ruta puramente técnica no ve: la percepción del declive de la lengua, la presión de las lenguas dominantes en la educación, la expectativa de ser incluidos en los sistemas de traducción existentes y la necesidad real de algo «no perfecto, pero suficientemente útil».

\lecturefigure{slide-07-speaker-interviews.jpg}{Alcance de las entrevistas a hablantes nativos de lenguas de bajos recursos}{Video de la clase de Stanford 00:05:54}

\begin{knowledgebox}{Lectura de la figura: las cifras de la muestra indican el alcance exploratorio, no permiten inferencias sobre la población}
Los 44 entrevistados y las 36 lenguas muestran que el equipo recogió opiniones de manera deliberada entre lenguas distintas, pero a cada lengua le corresponden muy pocas personas, y la propia clase reconoce que los canales de reclutamiento tienen sesgos regionales y hacia grupos migrantes. El uso correcto es formular design hypotheses y seguir consultando en las etapas de traducción, evaluación y publicación, no presentar las proporciones de las entrevistas como conclusiones estadísticas globales.
\end{knowledgebox}

\teachervoice{Fan llama a este paso un social-sciences-type o sociology-type approach. No afirma que las entrevistas puedan sustituir a los experimentos, sino que, en problemas cercanos a las personas, solo después de aclarar el objetivo se sabe qué deben medir los experimentos posteriores.}

\subsection{Declive de la lengua, inclusión tecnológica y lo «suficientemente bueno»}

Los resultados de las entrevistas situaron la traducción automática en un contexto social más amplio: a algunos usuarios les preocupa el declive de la lengua local en la educación y en la vida pública; al mismo tiempo, ser incluidos en los sistemas de traducción dominantes genera fuertes expectativas. La clase señala además que los usuarios no necesariamente exigen que la traducción automática cumpla todos los estándares de la traducción humana profesional; en ciertos escenarios de comunicación cotidiana, una herramienta con defectos evidentes pero que va en la dirección correcta también puede aportar valor.

\lecturefigure{slide-08-interview-findings.jpg}{Preocupaciones, expectativas y juicios de utilidad revelados por las entrevistas}{Video de la clase de Stanford 00:07:03}

\begin{knowledgebox}{Lectura de la figura: los tres hallazgos actúan sobre tres niveles distintos del sistema}
El declive de la lengua afecta las prioridades del proyecto y el riesgo cultural; la expectativa de «ser incluidos en el sistema» afecta los objetivos de cobertura; que la traducción automática pueda ser ya suficiente en el uso real exige que la evaluación distinga escenarios como la publicación profesional, la educación, los viajes y la salud. Un puntaje promedio no puede sustituir estas distintas tolerancias al riesgo.
\end{knowledgebox}

\begin{warningbox}{«Suficientemente útil» no significa rebajar el estándar para las lenguas de bajos recursos}
Los umbrales por escenario permiten que tareas distintas tengan tolerancias al error distintas, pero eso no justifica dejar a las lenguas de bajos recursos para siempre con una calidad inferior. Sobre todo en información médica, de salud pública, legal y de crisis, el error de traducción de una palabra pequeña puede tener consecuencias graves; el ejemplo de toxicity más adelante lo mostrará en concreto.
\end{warningbox}

\teachervoice{En la sesión de preguntas y respuestas de la clase se enfatiza de nuevo que la consulta con hablantes nativos debe abarcar todo el proceso de desarrollo: desde las entrevistas iniciales, pasando por el equipo de traducción profesional, hasta la evaluación humana y la retroalimentación abierta de la comunidad. Human-centered es una propiedad de todo el ciclo de vida, no un workshop de una sola vez.}

\subsection{Resumen de la sección}

Antes de modelar, NLLB confirmó las necesidades reales y extendió la participación de la comunidad a los datos y la evaluación. Las entrevistas aportan evidencia de diseño e indicios de riesgo, no estadísticas representativas; lo «útil» debe definirse por escenario y no puede convertirse en excusa para eludir la responsabilidad sobre la calidad.
\section{FLORES-200: primero, la infraestructura de evaluación}

Una vez establecida la necesidad real, el siguiente paso no es entrenar de inmediato, sino responder primero «cómo saber si el sistema de verdad mejora». Un problema frecuente en la investigación con pocos recursos es que faltan el conjunto de entrenamiento, el conjunto de validación y las métricas; sin un conjunto de evaluación común, comparable y traducido por personas, las mejoras del modelo no pueden verificarse entre idiomas. La serie FLORES se creó para cubrir esta carencia de infraestructura.

\subsection{El orden de investigación Evaluation-first}

Esta sección explica por qué la clase aborda primero la evaluation y después el training. La clase coloca de forma explícita los evaluation datasets antes que los training datasets. Este orden expresa una disciplina de ingeniería: si aún no se han definido las condiciones de éxito, la distribución de prueba y los tipos de falla, las iteraciones de entrenamiento solo optimizarán el objetivo sustituto que resulte más cómodo de medir.

\lecturefigure{slide-09-evaluation-datasets.jpg}{NLLB trata primero los datos de evaluación y después los datos de entrenamiento}{Video de la clase de Stanford 00:07:24}

\begin{knowledgebox}{Lectura de la figura: la portada de sección también expresa prioridades}
Esta slide casi no tiene detalles técnicos, pero muestra la lógica con la que se organiza el curso: el benchmark no es un accesorio que se agrega cuando el modelo ya está terminado, sino el sistema de coordenadas que orienta el entrenamiento. FLORES-200 define primero la comparabilidad entre idiomas; solo así el data mining, el MoE y el curriculum posteriores pueden juzgar aciertos y errores en la misma escala.
\end{knowledgebox}

\teachervoice{Fan dice que quiere «hablar primero de los evaluation data, porque hacer bien la evaluation es sumamente importante». Esta frase debe leerse como un método de investigación: cuando no se tiene un conjunto de prueba confiable, no hay que tomar la loss de entrenamiento ni los resultados de unos cuantos idiomas con muchos recursos como representativos del conjunto.}

\subsection{De FLORES a FLORES-200}

A continuación hay que entender cómo el benchmark se amplió de una generación a otra. FLORES se construyó en un principio en torno a la evaluación de unos pocos idiomas con pocos recursos; después se amplió a FLORES-101, y NLLB llevó la cobertura hasta FLORES-200. En clase se mencionó que el nombre proviene de Facebook Low Resource y que, aunque la empresa cambió su nombre a Meta, el proyecto conservó el nombre FLORES.

\lecturefigure{slide-10-flores-101.jpg}{FLORES-101 sentó el precedente para la evaluación a gran escala con pocos recursos}{Video de la clase de Stanford 00:08:12}

\begin{knowledgebox}{Lectura de la figura: el valor del benchmark proviene de oraciones fuente comunes y de la traducción humana}
FLORES hace que muchos idiomas compartan un mismo grupo de oraciones con correspondencia semántica, lo que permite comparar la traducción into-English, out-of-English y non-English. No consiste en recopilar páginas web existentes y dividirlas al azar, sino en producir puntos de referencia de evaluación más estables mediante traducción profesional y procesos de control de calidad.
\end{knowledgebox}

\begin{importantbox}{Por qué las direcciones de traducción crecen de forma combinatoria}
Si hay $N$ idiomas y se consideran $A\rightarrow B$ y $B\rightarrow A$ como direcciones distintas, existen como máximo
\[
D=N(N-1)
\]
direcciones de traducción dirigidas. $N$ representa el número de idiomas y $D$, el número de direcciones sin contar la copia al mismo idioma. Para $N=200$, $D=39{,}800$; esto explica las más de cuarenta mil direcciones que se mencionaron en clase, y también por qué es imposible que los datos de entrenamiento sean igual de abundantes para cada dirección.
\end{importantbox}

\subsection{Los cuatro atributos de diseño de FLORES}

Esta sección desglosa «tener un conjunto de evaluación» en decisiones de diseño concretas. La clase muestra las decisiones de diseño de FLORES con progressive bullets, cuyo estado final incluye cuatro puntos: centrarse en idiomas con pocos recursos, admitir many-to-many, abarcar temas y dominios diversos, y conservar el document-level context. Cada uno de estos atributos responde, respectivamente, al sesgo de cobertura, al sesgo de dirección, al sesgo de dominio y al sesgo de oraciones aisladas.

\lecturefigure{slide-11-flores-properties.jpg}{El diseño de FLORES-200: pocos recursos, muchos a muchos, varios dominios y nivel de documento}{Video de la clase de Stanford 00:09:57}

\begin{knowledgebox}{Lectura de la figura: cuatro atributos frente a cuatro distorsiones del benchmark}
Elegir solo idiomas con muchos recursos sobrestima la madurez del sistema; evaluar solo English-centric directions oculta la degradación en las direcciones non-English; evaluar un solo dominio hace pasar la coincidencia de dominio por una capacidad general; evaluar solo oraciones aisladas pasa por alto los pronombres, la coherencia y las ambigüedades que dependen del contexto. FLORES no puede eliminar todos los sesgos, pero reduce notablemente estas distorsiones y permite discutirlas.
\end{knowledgebox}

\begin{warningbox}{Many-to-many evaluation no equivale a many-to-many supervision}
Que el conjunto de evaluación permita construir pares de comparación para todos los idiomas no significa que en la etapa de entrenamiento se disponga de corpus paralelos para todas las direcciones. En la sesión de Q\&A de la clase se dijo explícitamente que NLLB no extrajo las $200\times200$ direcciones completas; muchas direcciones dependen de la zero-shot transfer que surge de una representación multilingüe compartida.
\end{warningbox}

\subsection{Traducción profesional, revisión automática y revisión humana}

El proceso de producción de FLORES comienza con el acuerdo sobre los estándares lingüísticos; después, los documentos se entregan a los traductores y pasan por revisiones automáticas, un reviewer humano y post-editing. Si la calidad queda por debajo del umbral, la muestra regresa a la etapa de corrección. Este ciclo muestra que los datos de evaluación de alta calidad no se obtienen con una sola compra, sino que requieren la acción conjunta de estándares lingüísticos, herramientas de revisión y juicio profesional.

\lecturefigure{slide-12-flores-collection-pipeline.jpg}{El proceso de FLORES, desde el acuerdo sobre estándares lingüísticos hasta la revisión y la posedición}{Video de la clase de Stanford 00:11:30}

\begin{knowledgebox}{Lectura de la figura: los ciclos entre flechas importan más que el nodo final de término}
Primero se fijan los estándares de script, ortografía y variante, para que los traductores sepan qué forma debe adoptar el texto de destino; las revisiones automáticas detectan omisiones de traducción, problemas de formato y caracteres anómalos, pero no pueden juzgar por sí solas la naturalidad semántica; el reviewer y el post-editing se encargan de regresar las muestras que no pasan. El `Quality > 90%` de la figura es la condición que se mostró en clase para este proceso y no debe extrapolarse como un umbral fijo de control para todos los proyectos lingüísticos.
\end{knowledgebox}

\teachervoice{El docente enfatiza que el primer paso, y el verdaderamente difícil, suele ser la alignment on language standards. Si una comunidad usa al mismo tiempo varios sistemas de escritura, ortografías y variantes locales, el equipo debe decidir primero qué se evalúa; no puede esperar a que termine el entrenamiento del modelo para completar las etiquetas.}

\subsection{Traductores, estándares, sistemas de escritura y variantes locales}

Por último, volvemos a las restricciones de producción del conjunto de evaluación. La última slide sobre datos de evaluación resume las dificultades de recopilación en tres tipos: encontrar traductores adecuados, coordinar los estándares lingüísticos y manejar la diversidad de sistemas de escritura y las variantes locales. Parecen asuntos de «operación de datos», pero en realidad determinan directamente si el benchmark mide el idioma de destino o solo una forma de escritura institucionalizada.

\lecturefigure{slide-13-flores-collection-challenges.jpg}{Los retos de traductores y estándares lingüísticos en la recopilación de datos de FLORES}{Video de la clase de Stanford 00:12:54}

\begin{knowledgebox}{Lectura de la figura: cada bullet se propaga al modelo y a las métricas}
Un reclutamiento insuficiente de traductores amplía la varianza; una mala elección de estándar hace que variantes legítimas se califiquen como errores; la mezcla de sistemas de escritura afecta la tokenización y la language identification; la ausencia de variantes locales hace que el modelo obtenga buenos resultados en el benchmark sin poder atender a las comunidades reales. El estándar de datos no es un recipiente neutral, sino el measurement model de todo el sistema.
\end{knowledgebox}

\begin{warningbox}{Translationese y English-centric source}
Un conjunto de evaluación traducido por personas aún puede contener translationese, es decir, oraciones de destino que conservan la estructura del idioma fuente; si el contenido fuente proviene sobre todo de un contexto en inglés, también puede carecer de temas culturales locales. La traducción humana de alta calidad mejora notablemente la comparabilidad, pero no equivale de forma automática a la distribución de un corpus nativo.
\end{warningbox}

\subsection{Resumen de la sección}

FLORES-200 establece primero un sistema de coordenadas de evaluación común para doscientos idiomas y, sobre esa base, hace comparable la investigación de modelos. Su valor proviene de la cobertura de idiomas con pocos recursos, el diseño muchos a muchos, la variedad de dominios, el contexto de documento y el proceso de traducción profesional; sus limitaciones incluyen la elección de estándares, el translationese y el sesgo del contenido fuente.
\section{De la semilla al corpus multilingüe explotable}

Un conjunto de evaluación responde «qué tan bien se hace»; los datos de entrenamiento además deben responder «de dónde se aprende». Para los idiomas de bajos recursos, la mayor dificultad no es la falta de un modelo lo bastante grande, sino que incluso falta un punto de partida confiable para language identification, vectores de oraciones, limpieza del corpus y backtranslation. NLLB-Seed y el mining pipeline posterior usan precisamente una pequeña cantidad de evidencia de alta calidad para poner en marcha una línea de producción de supervisión débil a mayor escala.

\subsection{Por qué no se puede empezar desde cero}

Esta sección pasa de la infraestructura de evaluación al problema de arranque de los datos de entrenamiento. La minería de datos suele describirse como «buscar oraciones paralelas en internet», pero para determinar a qué idioma pertenece una página web, si dos oraciones son semánticamente equivalentes o si un filtro descarta por error datos válidos, se necesitan anotaciones o representaciones confiables previas. La función de los seed data no es aportar la escala final, sino calibrar esos componentes básicos.

\lecturefigure{slide-14-seed-data-why.jpg}{La identificación de idioma, los vectores de oraciones y los modelos de traducción requieren datos semilla confiables}{Stanford, video de la clase 00:13:51}

\begin{knowledgebox}{Lectura de la figura: un mismo conjunto de seed data sirve a cuatro arrancadores}
Un texto confiable en el idioma de destino puede entrenar language identification; los pares de oraciones alineadas pueden entrenar un sentence encoder; el lado monolingüe puede apoyar el language modeling; el lado bilingüe puede poner en marcha un translation model. Si los datos iniciales están contaminados con idiomas equivocados, traducción automática o mezcla de escrituras, el error se propaga simultáneamente por las cuatro rutas.
\end{knowledgebox}

\begin{importantbox}{El efecto de palanca de los seed data}
El valor de una pequeña cantidad de datos humanos no depende solo del número de muestras, sino de cuántos componentes posteriores puede calibrar:
\[
\text{Seed}\rightarrow\{\text{LID},\text{Encoder},\text{Filter},\text{MT}\}
\rightarrow\text{Mined/BT Data}\rightarrow\text{Better MT}.
\]
Este es un bootstrapping loop. La calidad inicial determina el límite superior del error de las pseudoetiquetas y del corpus obtenido por minería en etapas posteriores; por ello, el valor marginal de la semilla puede ser mayor que el de la misma cantidad de texto web aleatorio.
\end{importantbox}

\subsection{NLLB-Seed: un recurso de arranque pequeño y confiable}

La slide de la clase informa que NLLB-Seed cubre 43 idiomas y alrededor de 6,193 oraciones, y explica que el contenido proviene de la lista de temas de cultura general de Wikimedia. Las cifras de la publicación final del artículo de NLLB varían ligeramente según la versión y la selección de idiomas; esta clase conserva las cifras de la slide y las interpreta como un recurso semilla de alta calidad de «unos miles de oraciones y decenas de idiomas», no como el grueso del entrenamiento a gran escala.

\lecturefigure{slide-15-nllb-seed.jpg}{Número de idiomas, número de oraciones y origen en temas de cultura general de NLLB-Seed}{Stanford, video de la clase 00:14:18}

\begin{knowledgebox}{Lectura de la figura: por qué se eligieron temas de cultura general y no páginas web aleatorias}
Los temas de cultura general abarcan personas, historia, geografía, ciencia y otros tipos de contenido, y reducen el ruido de plantillas de un solo sitio; además, las oraciones consecutivas dan contexto a los traductores. El costo es que el contenido de origen sigue la estructura de conocimiento de la Wikipedia en inglés y el texto de destino puede presentar translationese. El objetivo de la semilla es un arranque confiable, no representar todas las culturas locales ni los contextos de la lengua hablada.
\end{knowledgebox}

\teachervoice{Fan resume las tareas de generación de bajos recursos con «can't start from nothing»: sin unas cuantas oraciones confiables, es difícil saber en qué idioma está un texto y también construir la traducción inicial. El punto pedagógico de los seed data es la calidad del punto de partida, no una competencia por el número de muestras.}

\subsection{WikiMatrix: el antecedente de la minería a gran escala}

WikiMatrix mostró una posibilidad fundamental: usar vectores de oraciones multilingües para buscar oraciones semánticamente correspondientes en Wikipedia permite producir datos paralelos para un gran número de pares de idiomas. Para NLLB demostró que la embedding-based mining puede ampliar la cobertura, pero la escala y los temas de Wikipedia son limitados; para llegar realmente a doscientos idiomas hay que recurrir a la web abierta, mucho más desordenada.

\lecturefigure{slide-16-wikimatrix.jpg}{WikiMatrix usa vectores de oraciones multilingües para extraer oraciones paralelas a gran escala}{Stanford, video de la clase 00:15:00}

\begin{knowledgebox}{Lectura de la figura: los 135M y 1620 del título son escala, no garantía de calidad}
Cuantos más candidatos se extraen, más falsos positivos puede haber. La aportación de WikiMatrix fue convertir la recuperación por vectores de oraciones en un pipeline escalable; NLLB todavía necesitó un mejor encoder, identificación de idioma, margin score, filtrado y validación humana para convertir oraciones web similares en bitext utilizable para entrenamiento.
\end{knowledgebox}

\subsection{Sentence Alignment: de las páginas web a los pares de oraciones}

La clase divide la minería de pares de oraciones en tres niveles: buscar en Common Crawl contenido web posiblemente relacionado, extraer y limpiar el texto, y luego usar vectores de oraciones para encontrar coincidencias. La falla de cualquier nivel produce emparejamientos erróneos: las plantillas de las páginas pueden parecerse aunque el texto principal sea distinto, el segmentador de oraciones puede no adaptarse a la escritura y los vectores pueden acercar oraciones del mismo tema pero con hechos distintos.

\lecturefigure{slide-17-sentence-alignment.jpg}{Del emparejamiento de contenido de Common Crawl a la alignment a nivel de oración}{Stanford, video de la clase 00:15:39}

\begin{knowledgebox}{Lectura de la figura: la alignment no es coincidencia de cadenas}
El sistema primero reduce los documentos candidatos, luego codifica las oraciones en vectores compartidos y finalmente compara su cercanía semántica. Al leer la figura hay que distinguir content retrieval de sentence retrieval: la primera reduce el espacio de búsqueda; la segunda decide si se forma un ejemplo positivo. Si se aplica directamente el vecino más cercano global, las plantillas comunes y las oraciones cortas sepultan las traducciones verdaderas.
\end{knowledgebox}

Sean $f(x)$ y $g(y)$ los vectores de la oración de origen $x$ y de la oración de destino $y$, respectivamente; la similitud más básica es
\[
\operatorname{cos}(x,y)=\frac{f(x)^\top g(y)}{\lVert f(x)\rVert_2\lVert g(y)\rVert_2}.
\]
Aquí $f,g$ denotan el sentence encoder multilingüe; el numerador es el producto interno y el denominador normaliza los vectores. En la práctica, la mining suele usar un margin-based score, que compara la similitud del candidato con la similitud promedio de su vecindario local para reducir la interferencia de la hubness y de las oraciones genéricas.

\begin{lstlisting}[caption={Flujo conceptual de la minería de corpus paralelos con vectores de oraciones}]
for source_document in crawl:
    target_candidates = retrieve_related_documents(source_document)
    source_sentences = clean_and_segment(source_document)
    target_sentences = clean_and_segment(target_candidates)
    pairs = nearest_neighbors(source_sentences, target_sentences)
    keep = margin_filter(pairs, language_id=True, quality=True)
    send_selected_samples_to_bilingual_validation(keep)
\end{lstlisting}

\subsection{Encoder quality: el primer cuello de botella de los idiomas de bajos recursos}

Si el sentence encoder solo forma un buen espacio compartido para los idiomas de altos recursos, las oraciones de bajos recursos se agrupan con vecinos equivocados, y por grande que sea la mining posterior solo producirá ruido en masa. La clase compara masked language modeling con multilingual distillation: la primera aprovecha el contexto monolingüe; la segunda hace que un student multilingüe imite el espacio semántico de un teacher más fuerte.

\lecturefigure{slide-18-encoder-training.jpg}{Mejora del encoder con masked language modeling y multilingual distillation}{Stanford, video de la clase 00:16:33}

\begin{knowledgebox}{Lectura de la figura: los dos objetivos complementan, respectivamente, el modelado del lenguaje y la alineación entre idiomas}
El masked language modeling hace que el encoder aprenda la estructura interna del idioma de destino; la multilingual distillation usa oraciones semánticamente equivalentes para acercar distintos idiomas a un espacio compartido. La primera no garantiza por sí sola la alineación entre idiomas, y la segunda depende de un teacher confiable y de datos emparejados; NLLB combina ambas para aumentar el recall y la precision de la mining en idiomas de recursos extremadamente bajos.
\end{knowledgebox}

\begin{warningbox}{La puntuación del encoder se convierte en un sesgo de selección de datos}
El mining pipeline solo ve las similitudes que el encoder puede representar. Si cierta escritura, dominio o variante tiene mala calidad en el encoder, menos de sus oraciones entran al conjunto de entrenamiento y después al modelo le cuesta todavía más mejorar ese idioma, lo que forma una retroalimentación negativa del data flywheel.
\end{warningbox}

\subsection{LASER3: primero hacer más equitativa la representación, luego hablar de escala}

Ya se explicó que el encoder determina lo que la mining puede ver; esta sección verifica ese cuello de botella con los resultados de LASER3. La gráfica de resultados compara, por idioma, la calidad de los vectores de oraciones. Lo que la clase quiere que el lector observe no es que «todas las barras suben», sino que el nuevo encoder mejora notablemente muchos idiomas que antes eran muy débiles, lo que aumenta la probabilidad de que encuentren candidatos correctos en la minería web.

\lecturefigure{slide-19-laser3-encoder-quality.jpg}{Calidad del encoder interlingüe de LASER3 frente a las primeras versiones de LASER}{Stanford, video de la clase 00:18:48}

\begin{knowledgebox}{Lectura de la figura: primero buscar los idiomas con línea base baja y luego ver si la mejora es generalizada}
El eje horizontal son los idiomas y el eje vertical, una métrica de calidad relacionada con el encoder; las barras claras y oscuras corresponden a versiones distintas. Lo más importante no es el aumento promedio, sino si los idiomas de la cola pasan de ser casi inutilizables a poder sostener la mining. La figura tampoco demuestra que la calidad final de la traducción aumente en la misma proporción, porque después siguen los cuellos de botella de visibilidad web, filtrado y entrenamiento del modelo.
\end{knowledgebox}

\teachervoice{El docente llama a la encoder quality un low-resource challenge: un método de mining viable en idiomas de altos recursos, si la representación compartida falla para ciertos idiomas, los excluye antes de que lleguen al modelo.}

\subsection{Una fábrica de datos iterativa, no un ETL único}

La slide completa conecta en un ciclo language identification, monolingual cleaning, sentence mining, filtering, entrenamiento del modelo y bilingual validation. Cada mejora del modelo y del encoder puede redescubrir mejores datos; cada vez que la validación humana detecta un error sistemático, también hay que volver a la LID, a la limpieza o a la configuración de umbrales.

\lecturefigure{slide-20-data-pipeline.jpg}{Pipeline de datos de bajos recursos de NLLB y su ciclo de validación humana}{Stanford, video de la clase 00:21:15}

\begin{knowledgebox}{Lectura de la figura: seguir las flechas para encontrar la retroalimentación, no solo leer los recuadros}
Los LID training data sustentan la clasificación de idioma de las páginas web; los monolingual data limpios entran a la mining; el bitext candidato, tras el filtering, entrena el MT; las salidas del modelo y las muestras vuelven a ser validadas por personal bilingüe. Los múltiples ciclos de la figura significan que datos, modelo y evaluación iteran juntos, y que cualquier umbral fijo debe recalibrarse por idioma.
\end{knowledgebox}

\begin{importantbox}{Propagación del error en el pipeline}
Si la precision efectiva de la etapa $i$ es $p_i$ y las etapas se aproximan como filtros en serie, la confiabilidad aproximada de los candidatos finales está acotada por
\[
P_{\text{effective}}\approx\prod_i p_i
\]
Aquí $p_i$ representa la tasa de acierto de etapas como LID, limpieza, emparejamiento de documentos y puntuación de pares de oraciones. La expresión no es un modelo probabilístico exacto, pero muestra que varios componentes «bastante buenos» conectados en serie pueden perder calidad de forma significativa; por eso se necesitan validación humana por muestreo y corrección mediante ciclos.
\end{importantbox}

\subsection{Stopes: abrir el pipeline de investigación como infraestructura}

Esta sección pasa del flujo algorítmico a la ingeniería reproducible. A continuación, la clase presenta el punto de acceso de código abierto de Stopes. Lo que se abre no es solo el modelo final, sino también las herramientas para procesar páginas web monolingües, extraer bitext y reproducir la producción de datos. Este tipo de release permite que investigadores externos revisen la lógica de filtrado, sustituyan componentes de idioma y añadan datos locales, en lugar de limitarse a descargar un checkpoint inexplicable.

\lecturefigure{slide-21-stopes-open-source.jpg}{Stopes: la herramienta de minería de datos multilingües que NLLB publicó como código abierto}{Stanford, video de la clase 00:21:45}

\begin{knowledgebox}{Lectura de la figura: el nivel de apertura determina la profundidad de la reproducibilidad}
Abrir solo el modelo permite reproducir únicamente la inference; abrir el benchmark permite reproducir la evaluation; solo abrir las herramientas de LID, encoder, mining y filtering permite reconstruir los training data. Stopes está en el tercer nivel: no sustituye el conocimiento local de los idiomas, pero hace más verificables los supuestos de ingeniería del pipeline de datos.
\end{knowledgebox}

\subsection{El límite inferior estricto de los idiomas de recursos extremadamente bajos}

Por último, hay que delimitar el límite inferior estricto que el pipeline no puede superar y distinguir entre un recall algorítmico insuficiente y la inexistencia de suficiente texto en internet. Aun con mining a gran escala, algunos idiomas casi no tienen suficiente texto web monolingüe; una escritura particular puede hacer fallar la limpieza y la segmentación genéricas, y el contenido en línea puede concentrarse en lo religioso o en un solo dominio. En ese caso, «rastrear más páginas web» no produce automáticamente señales de entrenamiento diversas, equilibradas y confiables.

\lecturefigure{slide-22-very-low-resource-challenges.jpg}{Cuello de botella monolingüe, escritura y limitaciones de dominio en idiomas de recursos extremadamente bajos}{Stanford, video de la clase 00:22:36}

\begin{knowledgebox}{Lectura de la figura: tres desafíos corresponden a tres remedios no intercambiables}
La escasez de datos monolingües requiere datos de la comunidad, voz u otras fuentes; una escritura particular requiere normalization, segmentación y LID a la medida; la concentración en un dominio requiere ampliar deliberadamente los temas, no simplemente repetir el muestreo. Un modelo más grande no puede crear evidencia textual que no existe; solo ajusta con más detalle el sesgo existente.
\end{knowledgebox}

\begin{warningbox}{La escala no atraviesa el data floor}
Cuando la web presence se acerca a cero, el límite superior del recall de la mining lo determina la existencia de los datos; cuando las páginas web pertenecen a un solo dominio, aumentar la escala del crawl produce sobre todo repeticiones. Un plan de ingeniería debe diagnosticar por separado «faltan datos» y «el algoritmo no encontró los datos».
\end{warningbox}

\teachervoice{Fan afirma explícitamente que, para los idiomas de recursos extremadamente bajos, incluso la mining a gran escala es muy difícil, y que los monolingual data son el limiting factor. Este juicio reubica el problema: deja de ser la falta de capacidad de cómputo y pasa a ser la digitalización de los idiomas y la visibilidad de los datos.}

\subsection{Resumen de la sección}

NLLB usa una semilla de alta calidad para poner en marcha la LID, el encoder, la mining y el MT, y luego construye con LASER3 y Stopes una fábrica de datos iterativa. El límite superior del sistema está acotado a la vez por la calidad de la representación, la visibilidad web, el manejo de las escrituras, la diversidad de dominios y la validación humana.
\section{La receta de datos: de los atributos de cada fuente al equilibrio de la cola larga}

Contar con el mining pipeline no basta para mezclar directamente todos los pares de oraciones en el entrenamiento. Las distintas fuentes de datos difieren enormemente en alineación humana, ruido, escala y dependencia de modelos; estos atributos determinan la correlación de los errores, la escalabilidad y el sesgo. La clase primero enumera los retos del modelo, después usa un diagrama del Transformer para establecer el contexto mínimo y, por último, coloca los cinco tipos de datos en una misma tabla comparativa.

\subsection{Tres tipos de retos en la etapa del modelo}

Esta sección continúa con el pipeline de datos y explica por qué tener más datos se convierte de inmediato en un problema de modelado. Al crecer el volumen de datos, el modelo debe manejar a la vez el aumento efectivo de datos, la interferencia entre idiomas y la ampliación de capacidad. Los idiomas de bajos recursos necesitan compartir representaciones con los de altos recursos para obtener transfer, pero compartir en exceso puede producir interference; aumentar la capacidad alivia la competencia, pero hace que las tareas escasas se sobreajusten más rápido.

\lecturefigure{slide-23-model-challenges.jpg}{Retos acoplados del aumento de datos, la interferencia entre idiomas y el tamaño del modelo}{Video de la clase de Stanford 00:23:36}

\begin{knowledgebox}{Lectura de la figura: los tres bullets forman una tensión, no son TODO independientes}
Data augmentation amplía la cobertura, pero puede introducir ruido generado por modelos; reducir la language interference suele depender de más capacidad especializada, pero eleva el riesgo de sobreajuste; ampliar el modelo, a su vez, aumenta la complejidad del entrenamiento y del enrutamiento. Los MoE, EOM y el curriculum que se ven más adelante son precisamente una respuesta conjunta a este triángulo de contradicciones.
\end{knowledgebox}

\subsection{La ruta de traducción del Transformer encoder-decoder}

NLLB usa un encoder-decoder translation model. El encoder codifica la oración fuente completa en una representación contextualizada, y el decoder predice la oración objetivo token por token, condicionado por la representación de la fuente y por el prefijo objetivo ya generado. A diferencia del prompting con decoder-only, esta estructura separa explícitamente la comprensión del lado fuente y la generación del lado objetivo.

\lecturefigure{slide-24-transformer-preliminaries.jpg}{Flujo encoder-decoder del Transformer en la traducción automática}{Video de la clase de Stanford 00:24:03}

\begin{knowledgebox}{Lectura de la figura: las pilas superior e inferior asumen condiciones distintas}
El encoder de arriba lee simultáneamente los tokens fuente y forma representaciones entre posiciones; el decoder de abajo usa una causal mask, solo puede ver el prefijo objetivo y accede al encoder mediante cross-attention. En el entrenamiento suele usarse teacher forcing, mientras que en la inferencia el modelo depende de los tokens que él mismo generó antes, por lo que los errores pueden acumularse paso a paso.
\end{knowledgebox}

Dada una secuencia fuente $x$ y una secuencia objetivo $y=(y_1,\ldots,y_T)$, el objetivo de generación condicional es
\[
\mathcal{L}_{\text{MT}}(x,y)=-\sum_{t=1}^{T}\log p_\theta(y_t\mid y_{<t},x).
\]
$\theta$ representa los parámetros del modelo, $y_{<t}$ el prefijo objetivo, $p_\theta$ la probabilidad condicional del siguiente token y $T$ la longitud del objetivo. Todas las diferencias entre fuentes de datos terminan modificando la distribución de los pares $(x,y)$ que aparecen en esta pérdida.

\subsection{Los cinco tipos de fuentes de datos no pueden compararse solo por cantidad}

La tabla de la clase va agregando, uno por uno, NLLB-Seed, public bitext, mined bitext, multilingual MT backtranslation y bilingual SMT backtranslation. La tabla final pide al lector comparar cuatro atributos: si hay alineación humana, si hay ruido, si la escala es limitada y si depende de modelos existentes.

\lecturefigure{slide-25-data-source-comparison.jpg}{Provenance, ruido, escala y dependencia de las fuentes de datos de entrenamiento de NLLB}{Video de la clase de Stanford 00:26:03}

\begin{knowledgebox}{Lectura de la figura: leer columna por columna importa más que memorizar los nombres fila por fila}
Los datos human-aligned suelen tener alta calidad pero escala limitada; el public bitext puede venir de diversas fuentes y tiene calidad desigual; los mined data son escalables pero dependen del encoder y del filtrado; la MMT/SMT backtranslation puede aprovechar datos monolingües, pero hereda el sesgo del teacher model. Ninguna fuente cumple a la vez con alineación humana, ausencia de ruido, escala ilimitada e independencia de modelos.
\end{knowledgebox}

\begin{knowledgebox}{Términos clave: fuentes de datos de entrenamiento}
\begin{tabularx}{\textwidth}{@{}L{0.19\textwidth}L{0.25\textwidth}X@{}}
\toprule
Fuente & Problema que resuelve & Riesgos principales y relación con el curso \\
\midrule
NLLB-Seed & Aporta pares de oraciones confiables para el arranque & Escala pequeña, centrado en el inglés, puede contener translationese \\
Public bitext & Reutiliza recursos bilingües existentes & Fuentes heterogéneas, con límites de limpieza y de licencia distintos \\
Mined bitext & Amplía el corpus paralelo a partir de páginas web & Depende de LID, del encoder y de umbrales; contiene ruido por pares mal emparejados \\
MMT backtranslation & Genera oraciones fuente sintéticas con un teacher multilingüe & El sesgo del teacher se comparte entre idiomas; calidad inestable en bajos recursos \\
SMT backtranslation & Genera oraciones fuente sintéticas con un modelo estadístico bilingüe & Puede aportar un sesgo inductivo distinto, pero su cobertura y calidad son limitadas \\
\bottomrule
\end{tabularx}
\end{knowledgebox}

\subsection{Distribución original: unos cuantos idiomas concentran la gran mayoría de los datos}

La tabla de fuentes anterior explicaba la naturaleza de los datos; esta sección examina además la distribución de cantidades. Al combinar solo los public y seed data, el volumen de entrenamiento entre idiomas abarca varios órdenes de magnitud. En la figura, los idiomas de altos recursos forman barras altas y los de bajos recursos quedan pegados a la base; si se muestrea directamente según la frecuencia original, las actualizaciones del modelo quedan casi siempre dominadas por los idiomas grandes.

\lecturefigure{slide-26-public-seed-data.jpg}{Los public y seed data presentan una distribución de cola larga extrema}{Video de la clase de Stanford 00:26:33}

\begin{knowledgebox}{Lectura de la figura: el eje vertical está en escala logarítmica; las barras bajas no van ligeramente atrás}
Un eje vertical logarítmico implica que una diferencia de altura entre barras vecinas puede corresponder a diez o cien veces más datos. Primero hay que comparar la extensión de la cola larga y después observar idiomas de bajos recursos concretos. Un muestreo uniforme simple reutiliza pocos ejemplos y provoca sobreajuste; muestrear según la frecuencia natural hace que casi no aparezcan, por lo que se requiere la intervención conjunta del aumento de datos y de un muestreo controlado.
\end{knowledgebox}

\begin{warningbox}{Equilibrar los datos no es copiar cada idioma hasta la misma cantidad}
El muestreo repetido no agrega información nueva; al contrario, hace que las direcciones de bajos recursos memoricen más rápido el conjunto de entrenamiento. La estrategia de equilibrio debe considerar a la vez la diversidad efectiva de los ejemplos, el ruido, la dirección, la duración del entrenamiento y la capacidad del modelo.
\end{warningbox}

\subsection{Mining y backtranslation amplían la cola larga}

A continuación se observa cómo la augmentation modifica la cola larga. Al incorporar mined bitext y backtranslation, muchas direcciones de bajos recursos obtienen más volumen de entrenamiento y la cola larga se eleva. Pero la figura sigue mostrando diferencias marcadas, lo que indica que la función de la augmentation es mejorar la cobertura, no colocar a todos los idiomas en las mismas condiciones de datos.

\lecturefigure{slide-27-mining-backtranslation.jpg}{Mining y backtranslation amplían los datos de entrenamiento de bajos recursos}{Video de la clase de Stanford 00:27:00}

\begin{knowledgebox}{Lectura de la figura: cuando una barra crece, también hay que preguntar de dónde vienen los datos}
Los primary data son corpus alineados por personas o públicos; los mined data provienen de oraciones similares encontradas en la web; los BT data los genera un modelo, que produce oraciones fuente sintéticas a partir de oraciones objetivo monolingües. Al apilar los tres colores el total aumenta, pero su confiabilidad y su estructura de errores son distintas. Si un idioma de bajos recursos depende sobre todo de BT, el modelo puede reforzar en un ciclo cerrado los errores fijos del teacher.
\end{knowledgebox}

La backtranslation parte de una oración monolingüe $y$ en el idioma objetivo, usa un modelo inverso $q_\phi(x\mid y)$ para generar una oración fuente sintética $\hat{x}$ y luego entrena el modelo directo:
\[
\hat{x}\sim q_\phi(x\mid y),\qquad
\mathcal{L}_{\text{BT}}=-\log p_\theta(y\mid\hat{x}).
\]
$q_\phi$ es el teacher/backward model, $p_\theta$ es el modelo directo que se va a entrenar y $\hat{x}$ es la synthetic source. Su valor está en aprovechar grandes cantidades de texto monolingüe en el idioma objetivo; su riesgo es que los errores de $q_\phi$ se conviertan en etiquetas de entrenamiento.

\begin{lstlisting}[caption={Flujo conceptual de backtranslation con provenance y control de calidad}]
for target_sentence in target_monolingual_corpus:
    synthetic_source = backward_model.translate(target_sentence)
    score = quality_filter(synthetic_source, target_sentence)
    if score >= language_specific_threshold:
        training_buffer.add(
            source=synthetic_source,
            target=target_sentence,
            provenance="backtranslation",
            teacher=backward_model.version,
        )
\end{lstlisting}

\teachervoice{El giro central de la clase es que, después de «meter todos» los tipos de datos, el problema no termina. Más datos traen ruido y una cola larga más complejos, y el modelo necesita manejar la interferencia entre idiomas y el sobreajuste en bajos recursos.}

\subsection{Resumen de la sección}

La receta de datos de NLLB se compone del seed humano, el bitext público, el mined bitext y dos tipos de backtranslation. Cada fuente intercambia de manera distinta calidad, escala y dependencia de modelos; mejorar la cola larga no equivale a equilibrar la información, y el sobreajuste debe seguir controlándose desde el modelo y la programación del entrenamiento.
\section{MoE, sobreajuste y aprendizaje curricular}

Cuando doscientos idiomas y un gran número de direcciones comparten un modelo, la capacidad fija de una red dense debe atender al mismo tiempo tareas muy distintas entre sí. Mixture of Experts (MoE) amplía la capacidad de parámetros mediante enrutamiento condicional, pero los expertos dispersos también facilitan que las tareas de bajos recursos encuentren una ruta especializada capaz de memorizar el conjunto de entrenamiento. Por eso, la clase expone de manera consecutiva MoE, dropout, Expert Output Masking y curriculum learning.

\subsection{Token-level MoE: capacidad condicional, no capacidad gratuita}

MoE coloca varios expert en algunas capas feed-forward, y cada token se enruta solo a unos cuantos expertos. El número total de parámetros puede ser muy grande, mientras que el cómputo efectivo por token es disperso. Distintos idiomas o estructuras pueden desarrollar cierta specialization, lo que reduce la interference que surge cuando todas las tareas se concentran en la misma subcapa dense.

\lecturefigure{slide-28-moe-architecture.jpg}{Compromiso entre capacidad, interferencia y sobreajuste en Token-level Mixture of Experts}{Video de la clase de Stanford 00:27:45}

\begin{knowledgebox}{Lectura de la figura: las tres bullet forman una cadena causal}
Token-level MoE aumenta la capacidad condicional, lo que puede permitir que coexistan la compartición entre idiomas y la especialización de los expertos; las rutas de expertos pueden reducir parte de la interference; pero, cuando los token de bajos recursos entran una y otra vez en unos pocos expertos, la capacidad adicional memoriza más rápido los datos escasos. MoE resuelve la competencia por capacidad, pero convierte la regularization en un problema más central.
\end{knowledgebox}

Sea $h$ la representación de un token, con $E$ expertos $f_e$; el enrutador asigna los pesos $g_e(h)$ y el conjunto de los expertos Top-$k$ es $\mathcal{T}_k(h)$. Entonces, la salida de MoE puede escribirse como
\[
\operatorname{MoE}(h)=\sum_{e\in\mathcal{T}_k(h)} g_e(h)f_e(h).
\]
$E$ denota el número total de expertos, $k$ el número de expertos que activa cada token, $f_e$ es una expert FFN y $g_e(h)$ es el peso de enrutamiento normalizado. El total de parámetros crece con $E$, pero el cómputo por token crece principalmente con $k$.

\begin{knowledgebox}{Términos clave: componentes de MoE}
\begin{tabularx}{\textwidth}{@{}L{0.21\textwidth}L{0.26\textwidth}X@{}}
\toprule
Término & Mecanismo & Riesgo / relación con la clase \\
\midrule
Expert & Subred FFN independiente & Aporta capacidad condicional, pero también puede memorizar tareas pequeñas \\
Router / gate & Elige expertos para cada token & El sesgo de enrutamiento puede concentrar la carga o fijar la división del trabajo por idioma \\
Top-$k$ routing & Solo activa unos cuantos expertos & Controla el cómputo, pero requiere load balancing \\
Language interference & Las actualizaciones de una tarea perjudican a otras & Con demasiada compartición, los idiomas de bajos recursos quedan sepultados por los gradientes de los de altos recursos \\
Overfitting & El entrenamiento continúa y la validación empeora & Aparece antes en las direcciones de bajos recursos; requiere regularización y programación \\
\bottomrule
\end{tabularx}
\end{knowledgebox}

\subsection{Dense dropout, MoE dropout y EOM}

La figura triple de la clase compara un dense model, MoE con dropout convencional y MoE con Expert Output Masking (EOM). El dropout convencional es eficaz en el modelo dense, pero no basta para contener por completo el sobreajuste de MoE en las direcciones de bajos recursos; EOM enmascara aleatoriamente el expert output de una parte de los token, de modo que el modelo no pueda depender en cada actualización de la capacidad especializada que obtiene por enrutamiento.

\lecturefigure{slide-29-dropout-eom-results.jpg}{Curvas de validación de Dense, MoE, dropout y Expert Output Masking}{Video de la clase de Stanford 00:29:18}

\begin{knowledgebox}{Lectura de la figura: fila superior, bajos recursos; fila inferior, altos recursos; primero observa el punto de inflexión de las curvas}
El eje horizontal son las updates de entrenamiento y el eje vertical muestra la calidad o la perplexity de validación. Las direcciones de bajos recursos alcanzan antes el punto de inflexión en el que «el entrenamiento continúa pero la validación empeora»; las de altos recursos suelen seguir beneficiándose. El objetivo de EOM no es acelerar el entrenamiento, sino retrasar el sobreajuste de bajos recursos y conservar, en la medida de lo posible, la ventaja de capacidad de MoE.
\end{knowledgebox}

Para el token $t$ y el experto seleccionado $e$, EOM puede abstraerse como
\[
\widetilde{f}_e(h_t)=m_{t,e}f_e(h_t),\qquad
m_{t,e}\sim\operatorname{Bernoulli}(1-p_{\text{eom}}).
\]
$m_{t,e}$ es una mask aleatoria, $p_{\text{eom}}$ es la probabilidad de enmascaramiento y $f_e(h_t)$ es la salida del experto. La implementación real se combina con Top-$k$ de forma más detallada, pero lo esencial es omitir aleatoriamente una parte de la expert contribution, en lugar de aplicar solo dropout convencional a la dense activation final.

\begin{warningbox}{EOM no consiste en «apagar expertos defectuosos»}
La mask es una regularización aleatoria durante el entrenamiento; no desactiva de forma permanente a un experto según la calidad de un idioma. Reduce la dependencia del modelo respecto de rutas token-expert fijas; si se interpreta como expert pruning estático, se malinterpretan las curvas y el comportamiento en despliegue.
\end{warningbox}

\teachervoice{Fan describe la capacidad adicional de MoE con la expresión very easy to overfit. La clase no presenta MoE como una arquitectura superior a dense de forma incondicional, sino que la considera, junto con la regularization de bajos recursos, parte de un mismo diseño.}

\subsection{Curriculum learning: incorporar al entrenamiento según el momento de sobreajuste}

Esta sección pasa de la regularización espacial a la programación temporal. Los pares de idiomas difieren en volumen de datos y en el tiempo que pueden entrenarse. NLLB observa primero en qué momento empieza a sobreajustarse cada dirección en un entrenamiento sin curriculum, y después agrupa las direcciones por intervalos; cuanto antes se sobreajusta una dirección, más tarde se incorpora al entrenamiento, de modo que recibe menos updates, mientras que las direcciones de altos recursos entran antes y se entrenan durante más tiempo.

\lecturefigure{slide-30-curriculum-learning.jpg}{Agrupación por intervalos según el momento de sobreajuste de cada par de idiomas e incorporación al entrenamiento por fases}{Video de la clase de Stanford 00:30:21}

\begin{knowledgebox}{Lectura de la figura: el schedule se basa en el comportamiento en validación, no solo en el número de ejemplos}
Los pares de idiomas con pocos datos suelen sobreajustarse antes, pero con el mismo número de ejemplos el comportamiento puede variar por el ruido, la escritura o el efecto de transferencia. La propuesta de la clase mide primero el overfitting onset y luego define el bucket y el introduction time. Así, el curriculum se guía por un diagnóstico empírico, y no por un orden fijo de «primero altos recursos, después bajos recursos».
\end{knowledgebox}

Sea $T$ el número total de pasos de entrenamiento y supongamos que el par de idiomas $d$ empieza a sobreajustarse tras unos $k_d$ pasos en el entrenamiento base; entonces puede incorporarse en el momento
\[
t_{\text{start}}(d)=T-k_d
\]
$t_{\text{start}}(d)$ es el momento de incorporación, $T$ es el total de pasos y $k_d$ es la ventana de entrenamiento que esa dirección tolera. El sistema real usa la mediana de cada intervalo en lugar de un momento independiente por dirección, para reducir la complejidad de la programación.

\begin{lstlisting}[caption={Phased curriculum basado en el momento de sobreajuste en el conjunto de validación}]
baseline = train_without_curriculum(all_directions)
for direction in all_directions:
    onset[direction] = first_overfitting_step(baseline.validation(direction))

buckets = cluster_by_onset(onset)
model = initialize_model()
for step in range(total_steps):
    active = directions_whose_bucket_has_started(step, buckets)
    model.update(sample_batch(active))
\end{lstlisting}

\begin{warningbox}{Sumar curriculum y EOM no garantiza un mejor resultado}
La ablation del artículo de NLLB muestra que, en configuraciones más pequeñas donde EOM ya mitiga el sobreajuste, añadir curriculum puede no aportar nada o incluso empeorar ligeramente; en la configuración completa a gran escala, la combinación vuelve a mejorar las direcciones de bajos recursos. La conclusión debe ligarse al modelo, los datos y la escala de entrenamiento; no puede presentarse como una receta universal.
\end{warningbox}

\subsection{Resumen de la sección}

MoE aumenta la capacidad condicional y mitiga la interferencia mediante token-level routing, pero amplifica el sobreajuste en bajos recursos. EOM aplica regularización aleatoria al expert output, y el curriculum controla la ventana de entrenamiento según el momento de sobreajuste en el conjunto de validación; el efecto de ambos depende de la escala y de las condiciones de los datos concretas.
\section{Los resultados deben leerse desglosados por dirección, recursos y métrica}

Una vez entrenado el modelo, lo más riesgoso es reportar solo una puntuación promedio. Los doscientos idiomas incluyen into-English, out-of-English y non-English directions, así como high, low y very-low resource groups; el promedio permite que los grupos más numerosos o con puntuaciones más altas oculten las direcciones realmente débiles. En la clase se usan tres tablas para ejercitar paso a paso esta lectura desglosada.

\subsection{Primero, comparar con los sistemas multilingües de la generación anterior}

Esta sección pasa del mecanismo de entrenamiento a los resultados a nivel de sistema y comienza por establecer una línea base histórica. La tabla comparativa de FLORES-101 coloca a NLLB junto a sistemas anteriores como M2M-100 y DeepNet en el mismo evaluation setup. Al leer la tabla, conviene confirmar si coinciden el número de idiomas que cubre cada modelo, los parámetros y datos de entrenamiento, y la metric; solo después se observa el cambio en el promedio.

\lecturefigure{slide-31-flores101-comparison.jpg}{Comparación de NLLB con sistemas multilingües anteriores en FLORES-101}{Video de la clase de Stanford 00:32:06}

\begin{knowledgebox}{Lectura de la figura: primero las definiciones de columna y los faltantes, después el ganador}
La tabla presenta por separado los resultados into-English, out-of-English, non-English y el promedio; `-/-` indica que algunas fuentes no tienen resultados directamente comparables. La mejora de NLLB respalda la idea de una «mejora conjunta de datos, modelo y entrenamiento», pero una sola tabla no basta para aislar la contribución independiente de cada componente.
\end{knowledgebox}

\begin{warningbox}{La comparación entre generaciones de modelos debe controlar el evaluation protocol}
El tokenizador, el conjunto de idiomas, el checkpoint, la configuración de decodificación y la metric implementation alteran la puntuación. Si los protocolos difieren, los valores de la tabla solo sirven como referencia a nivel de sistema en el momento de la clase, y no deben reordenarse como un leaderboard permanente fuera de su contexto.
\end{warningbox}

\subsection{FLORES-200: desglose por nivel de recursos}

La tabla de FLORES-200 divide además out-of-English e into-English en all, high, low y very-low, y presenta non-English en una columna aparte. Así el lector puede ver si la mejora global llega de verdad a la cola de bajos recursos, en lugar de provenir solo de los idiomas con muchos recursos.

\lecturefigure{slide-32-flores200-results.jpg}{Resultados de FLORES-200 reportados por dirección y nivel de recursos}{Video de la clase de Stanford 00:33:09}

\begin{knowledgebox}{Lectura de la figura: primero comparar horizontalmente los niveles de recursos, luego verticalmente las direcciones}
En una misma fila, la caída de high a very-low muestra el efecto de la disponibilidad de recursos; dentro de un mismo nivel de recursos, la diferencia entre into y out-of-English muestra la asimetría de dirección; el agregado non-English depende de la transfer multilingüe. El promedio resume el sistema, pero no sustituye estos tres diagnósticos.
\end{knowledgebox}

En los resultados de la clase aparecen tanto BLEU como chrF++. BLEU compara principalmente la n-gram precision tokenizada y aplica una penalización por longitud; chrF++ se apoya más en n-gramas de caracteres e incorpora word n-grams, por lo que suele ser más robusta en idiomas morfológicamente ricos o con tokenización inestable. Ambas son solo reference-based automatic metrics.

\begin{knowledgebox}{Términos clave: las métricas de las tablas de resultados}
\begin{tabularx}{\textwidth}{@{}L{0.2\textwidth}L{0.27\textwidth}X@{}}
\toprule
Métrica/grupo & Qué problema resuelve & Qué no puede demostrar \\
\midrule
BLEU / spBLEU & Compara coincidencias de n-gramas con una tokenización unificada & No representa directamente el significado, la naturalidad ni la seguridad \\
chrF++ & Reduce la sensibilidad a la tokenización mediante n-gramas de caracteres y de palabras & Sigue dependiendo de la traducción de referencia y no cubre todas las expresiones válidas \\
Into-English & Entrada en múltiples idiomas hacia el inglés & No representa la calidad de generación cuando el destino es un idioma de bajos recursos \\
Out-of-English & Entrada en inglés hacia múltiples idiomas & Se ve afectada fácilmente por la generación en el lado de destino y la escasez de datos \\
Non-English & Traducción entre dos idiomas distintos del inglés & Suele depender de la transfer; la cobertura de direcciones en el entrenamiento es desigual \\
\bottomrule
\end{tabularx}
\end{knowledgebox}

\subsection{Los límites correctos de la comparación con sistemas comerciales}

En la clase se colocan además las direcciones de bajos recursos junto a sistemas externos como Google Translate. Esta comparación sirve para mostrar la competitividad de NLLB en cobertura de investigación y como modelo público, pero los productos externos se actualizan constantemente y sus idiomas soportados, su decodificación y sus datos internos no son visibles, por lo que no constituye una referencia estable.

\lecturefigure{slide-33-external-system-comparison.jpg}{Comparación de resultados entre NLLB y los sistemas comerciales de traducción en el momento de la clase}{Video de la clase de Stanford 00:33:27}

\begin{knowledgebox}{Lectura de la figura: una comparación con productos es, ante todo, una instantánea con fecha}
Primero hay que confirmar qué direcciones de idioma soportan ambos sistemas y después comparar las columnas de bajos recursos; un hueco no debe interpretarse como una puntuación de cero. Los sistemas comerciales pueden usar datos no publicados y reglas manuales, mientras que NLLB pone el énfasis en activos de investigación abiertos. La tabla respalda que NLLB «alcanza una calidad competitiva en algunas direcciones», no que «supere al producto en todos los escenarios».
\end{knowledgebox}

\begin{importantbox}{El 44\% en BLEU es una mejora relativa}
Si la puntuación anterior es $B_{\text{old}}$ y la nueva es $B_{\text{new}}$, la mejora relativa se define como
\[
r=\frac{B_{\text{new}}-B_{\text{old}}}{B_{\text{old}}}\times100\%.
\]
$r=44\%$ no significa que BLEU haya aumentado 44 puntos absolutos ni que cada par de idiomas haya mejorado 44\%. Es la mejora relativa global que reporta el proyecto NLLB bajo el conjunto de comparación y el protocolo especificados.
\end{importantbox}

\teachervoice{Fan ofrece un criterio práctico de la clase: con algunas métricas, alrededor de 30 el sistema empieza a ser «bastante utilizable». Ella lo llama rule of thumb, no un umbral universal; el idioma, la dirección, el dominio y el costo de los errores modifican la utilidad.}

\subsection{Resumen de la sección}

Los resultados deben desglosarse por translation direction, resource level y metric. Las tablas automáticas permiten comparar sistemas con rapidez, pero no demuestran la calidad percibida por el usuario, la seguridad ni una mejora uniforme en todas las direcciones; el 44\% es una mejora relativa sujeta a condiciones, no una garantía de calidad absoluta.
\section{Evaluación humana y seguridad: los riesgos reales más allá de la calidad promedio}

Las métricas automáticas son adecuadas para iterar con un gran número de experimentos, pero el destinatario final de un sistema multilingüe son las personas. Los evaluadores de distintos idiomas pueden usar escalas diferentes, y un reference-based score tampoco puede juzgar todos los problemas semánticos, de cortesía y culturales. Por eso, en clase se separan la human evaluation y la toxicity en flujos de trabajo independientes.

\subsection{Las automatic metrics son rápidas; la human evaluation se acerca más al uso real}

Esta sección plantea además si los resultados automáticos pueden representar la experiencia humana. La slide de la clase resume con un contraste la división de funciones entre ambas: las métricas automáticas son rápidas y adecuadas para la iteración de investigación; la evaluación humana es cara y lenta, pero se acerca más a la calidad real. No se trata de elegir una u otra, sino de usar las métricas automáticas para filtrar alternativas y el juicio humano para calibrar la conclusión final.

\lecturefigure{slide-34-human-evaluation.jpg}{Las métricas automáticas sirven para iterar rápido; la evaluación humana, para juzgar la calidad real}{Video de la clase de Stanford 00:33:45}

\begin{knowledgebox}{Lectura de la figura: la velocidad y la validez son dos ejes}
Las automatic metrics pueden calcularse rápidamente sobre cuarenta mil direcciones y permiten hacer ablation; la human evaluation puede juzgar la fidelidad semántica, la naturalidad y la aceptabilidad, pero está limitada por el costo, los criterios de evaluación y la cobertura de idiomas. Un proceso maduro necesita que las métricas automáticas aporten breadth y que la evaluación humana aporte validity.
\end{knowledgebox}

\begin{warningbox}{Un pequeño aumento en la puntuación automática no siempre es perceptible}
Un aumento de BLEU/chrF++ puede deberse a la morfología de las palabras o a la coincidencia con la referencia, y no mejora necesariamente la experiencia del usuario. Sobre todo cuando las puntuaciones son cercanas, hay que revisar los intervalos de confianza, la agrupación por idioma y las preferencias humanas, en lugar de presentar diferencias en los decimales como un avance de nivel de producto.
\end{warningbox}

\teachervoice{Fan llama directamente a la human evaluation the real deal. No descarta las automatic metrics, sino que advierte: el objetivo sustituto de la iteration de investigación no puede convertirse automáticamente en una conclusión sobre el uso real.}

\subsection{Las escalas de evaluación entre idiomas necesitan calibrarse}

El problema que resuelve el estudio `Consistent Human Evaluation' es cómo logran los evaluadores de distintos pares de idiomas dar puntuaciones comparables. Si la comunidad de algún idioma es en general más estricta, un promedio simple confundirá la diferencia de escala de los evaluadores con una diferencia entre modelos; por eso se necesitan instrucciones comunes para la tarea, capacitación, controles de calidad y una normalización adecuada.

\lecturefigure{slide-35-human-evaluation-paper.jpg}{Evaluación humana consistente de la traducción automática entre pares de idiomas}{Video de la clase de Stanford 00:34:27}

\begin{knowledgebox}{Lectura de la figura: el título del artículo enfatiza consistent, no solo human}
La evaluación humana no es objetiva por naturaleza. Los evaluadores necesitan comprender la oración de origen y la de destino, distinguir adequacy de fluency y seguir una rubric común. La consistencia entre idiomas determina si las distintas direcciones pueden compararse en una misma figura; de lo contrario, las diferencias de altura de las barras entre idiomas pueden mezclar diferencias de cultura y de escala de evaluación.
\end{knowledgebox}

Sea $n_d$ el número de puntuaciones humanas válidas $r_{d,i}$ de la dirección $d$; el promedio más simple es
\[
\bar r_d=\frac{1}{n_d}\sum_{i=1}^{n_d}r_{d,i}.
\]
$r_{d,i}$ es la $i$-ésima puntuación, $n_d$ es el tamaño de muestra válido de esa dirección y $\bar r_d$ es el promedio de la dirección. Una comparación real requiere además corregir por la incertidumbre, el efecto del evaluador y la dificultad de las muestras; reportar solo $\bar r_d$ oculta la varianza.

\subsection{Cómo leer los resultados humanos: into, out-of y non-English}

Una vez calibradas las escalas de evaluación, esta sección lee los resultados humanos finales. La gráfica de barras de la clase reúne los tres tipos de dirección. El progressive state final muestra la distribución de la calidad humana por idioma o grupo de idiomas, con el propósito de comprobar si las mejoras en las métricas automáticas se traducen en mejoras perceptibles para las personas y qué direcciones siguen notablemente rezagadas.

\lecturefigure{slide-36-human-evaluation-results.jpg}{Resultados de la evaluación humana de NLLB en tres tipos de dirección de traducción}{Video de la clase de Stanford 00:36:03}

\begin{knowledgebox}{Lectura de la figura: primero las diferencias dentro de cada grupo, luego los promedios entre grupos}
Las barras verdes, azules, etc., corresponden a los grupos into-English, non-English y out-of-English. Primero se buscan los valles dentro de cada grupo para determinar si las debilidades se concentran en la generación hacia idiomas de destino de bajos recursos; después se compara la posición general. La figura respalda que NLLB obtiene mejoras humanas amplias, pero también muestra que las diferencias entre idiomas están lejos de desaparecer.
\end{knowledgebox}

\begin{warningbox}{Una gráfica de barras no demuestra causalidad}
Los resultados humanos provienen del sistema completo y no pueden atribuirse por separado a LASER3, mining, MoE, EOM o curriculum. Para juzgar la contribución de cada componente hay que combinarlos con una controlled ablation; la figura del sistema final solo responde si la solución conjunta es eficaz sobre la distribución de evaluación.
\end{warningbox}

\subsection{Toxicity: no todos los errores de traducción son igual de graves}

Esta sección pasa de la calidad promedio a los riesgos de seguridad de la cola. La puntuación promedio de calidad puede ocultar errores catastróficos. En clase se da el ejemplo, en el contexto de COVID, de `wash hands' traducido como `hold hands': las dos frases se parecen en la superficie, pero su significado como acción es opuesto. En la información de salud pública, el costo de este tipo de error es mucho mayor que el de un orden de palabras ligeramente forzado.

\lecturefigure{slide-37-toxicity-error-severity.jpg}{Un ejemplo de salud pública muestra que los errores de traducción tienen distinta gravedad}{Video de la clase de Stanford 00:37:12}

\begin{knowledgebox}{Lectura de la figura: la gravedad del error depende del contexto y de las consecuencias de la acción}
Una metric común puede tratar la diferencia de una palabra como un mismatch local, pero el usuario actuará de manera completamente distinta a partir de ella. La evaluación de seguridad necesita identificar patrones de alto costo como meaning inversion, added toxicity, negaciones omitidas y ataques a la identidad, y reportarlos por separado según el dominio, en lugar de depender solo del BLEU promedio.
\end{knowledgebox}

\begin{importantbox}{La calidad y la seguridad deben reportarse en dos ejes}
La aceptación de la dirección $d$ puede escribirse como una condición bidimensional:
\[
\text{Accept}(d)=\mathbf{1}[q_d\ge\tau_q]\cdot
\mathbf{1}[a_d\le\tau_a],
\]
donde $q_d$ es la calidad general, $a_d$ es la tasa de added-toxicity u otros errores de alto riesgo, $\tau_q$ es el límite inferior de calidad y $\tau_a$ es el límite superior de riesgo. En clase no se fija un umbral único; la fórmula sirve para mostrar que una calidad promedio alta no compensa una falla de seguridad.
\end{importantbox}

\teachervoice{Fan pregunta «¿por qué me importa tanto la toxicity?» y luego explica con la mala traducción de salud pública que los errores de un sistema de traducción no son token mismatch equivalentes, sino que cambian conductas reales. Los requisitos de seguridad deben formar parte de la evaluación previa a la publicación del modelo.}

\subsection{Toxicity-200: recursos de detección ligados a la cultura}

Con ejemplos de errores de alto riesgo, el siguiente paso es construir recursos de detección escalables. El equipo recopiló toxicity lists para doscientos idiomas y subraya que estas listas de palabras son específicas de cada cultura. Una misma palabra puede tener distinto grado de ofensa, contexto y tabú en comunidades diferentes; traducir directamente una lista de groserías en inglés produce a la vez falsos negativos y falsos positivos.

\lecturefigure{slide-38-toxicity-lists.jpg}{Toxicity lists específicas de cada cultura para doscientos idiomas, usadas para detección y mitigación}{Video de la clase de Stanford 00:38:15}

\begin{knowledgebox}{Lectura de la figura: la lista de palabras es una sonda, no un modelo de seguridad completo}
Una toxicity list puede detectar palabras dañinas que no están en la oración de origen pero que la traducción añade, y es adecuada para un screening a gran escala; sin embargo, no puede entender citas, negaciones, términos reapropiados ni el contexto, y también puede pasar por alto expresiones ofensivas que no están en la lista. Su papel correcto es el de risk probe, en combinación con la revisión humana y la evaluación semántica.
\end{knowledgebox}

\begin{warningbox}{No hay que trasladar directamente una safety taxonomy del inglés}
La cultura, las categorías de identidad y las formas de agresión cambian de un idioma a otro. La recopilación de las listas debe calibrarse con hablantes nativos y conocimiento de la comunidad, y registrar la versión y el ámbito de aplicación; de lo contrario, un «estándar de seguridad unificado» puede clasificar como dañinas expresiones locales normales o pasar por alto riesgos reales.
\end{warningbox}

\subsection{Resumen de la sección}

Las métricas automáticas aportan breadth para la iteración, la evaluación humana aporta validity sobre la calidad real, y la toxicity revisa los errores de alto costo que la puntuación promedio no puede expresar. Las tres deben reportarse en paralelo, y tanto las escalas humanas como las listas de palabras culturales requieren una calibración localizada por idioma.
\section{Código abierto, direcciones futuras y preguntas y respuestas de la clase}

Al cierre de la clase, NLLB no se presenta como un punto final: se proponen tres direcciones futuras y, en la sesión de Q\&A, se añaden precisiones sobre la participación de la comunidad, el volumen del trabajo con datos, Common Crawl, zero-shot, speech translation y los límites de los foundation models. Este contenido oral explica cómo el sistema especializado NLLB de 2022 se conecta con los modelos multimodales de propósito general que llegaron después.

\subsection{Explicit multilinguality, support y facilidad de uso}

La slide de direcciones futuras destaca explicit multilinguality, support for everyone, y que el modelo sea más fácil de usar y de entrenar. El primer punto exige que el modelo conozca de forma explícita el idioma, la escritura y la dirección de traducción; el segundo incluye a los idiomas de la cola entre los objetivos; el tercero reconoce que una cadena de procesamiento de datos que solo pueden operar las grandes instituciones todavía no basta para constituir una infraestructura de acceso universal.

\lecturefigure{slide-39-future-directions.jpg}{Direcciones futuras de NLLB: multilingüismo explícito, cobertura más amplia y menores barreras de uso}{Video de la clase de Stanford 00:41:15}

\begin{knowledgebox}{Lectura de la figura: las direcciones futuras corresponden a tres niveles, representación, cobertura e ingeniería}
Explicit multilinguality actúa sobre los language tags, el routing y la evaluation; support for everyone actúa sobre los idiomas de la cola, las variantes dialectales y las necesidades de las comunidades; ease of use/training actúa sobre las herramientas abiertas, la barrera de cómputo y la adaptación local. Si falta cualquiera de los tres, los resultados de investigación difícilmente se convierten en una capacidad pública utilizable de forma sostenida.
\end{knowledgebox}

\teachervoice{Fan señala en particular que algunos idiomas son sobre todo orales y no escritos. Un NLLB de texto no puede cubrir esos idiomas a partir de nada; en el futuro habrá que conectar speech, transcription y text translation, en lugar de tratar la «ausencia de datos web» como si el idioma no existiera.}

\subsection{El trabajo con datos y el trabajo con modelos pesan casi lo mismo}

En la sesión de Q\&A alguien preguntó cómo se repartían los recursos del proyecto. Fan respondió con una estimación aproximada: cerca de la mitad fue trabajo de data o data-adjacent, y la otra mitad de modeling, aunque la frontera es muy difusa. LID, el encoder, el filtering y la evaluation pertenecen a la vez a los datos y al modelo; lo importante de esta respuesta no es una proporción exacta de gestión de proyecto, sino refutar la idea de que «el trabajo principal es entrenar un modelo grande».

\begin{importantbox}{Niveles de los activos de ingeniería de NLLB}
\begin{tabularx}{\textwidth}{@{}L{0.2\textwidth}L{0.29\textwidth}X@{}}
\toprule
Nivel & Activos públicos & Pregunta que permiten reproducir \\
\midrule
Need / standards & Conclusiones de las entrevistas y proceso de normalización lingüística & Por qué se definieron así los objetivos y las etiquetas \\
Evaluation & FLORES-200, human protocol, Toxicity-200 & Si el modelo mejora y dónde no es seguro \\
Data tools & LID, LASER3, Stopes, scripts de reconstrucción & Cómo se generan y filtran los datos de entrenamiento \\
Models & Checkpoints y código de NLLB & Comportamiento final en inferencia y adaptación \\
\bottomrule
\end{tabularx}
Una apertura completa significa cubrir, en la medida de lo posible, toda la cadena que va del problema al modelo, y no solo publicar el checkpoint.
\end{importantbox}

\teachervoice{Al decir «50-50», el docente añadió de inmediato que data y modeling difícilmente se separan con rigor. Estas notas lo toman como una intuición de orden de magnitud: en los grandes proyectos multilingües, una parte considerable de la innovación principal ocurre en los datos, la evaluación y la cadena de herramientas.}

\subsection{Que Common Crawl sea abierto no significa que represente por igual}

En clase se explicó que Common Crawl es un conjunto de instantáneas web abiertas y descargables, y se usó una formulación prudente sobre su frecuencia de actualización. El acceso abierto resuelve parte del problema de los permisos para obtener datos, pero no resuelve la proporción de páginas web por idioma, la duplicación, las plantillas, el contenido basura, el sesgo geográfico ni la deriva temporal.

\begin{warningbox}{Web-scale y language coverage son dos variables distintas}
Un corpus puede contener miles de millones de páginas web y, al mismo tiempo, casi no tener texto útil en ciertos idiomas. Al evaluar un mining pipeline conviene reportar, por idioma, las raw pages, la LID pass rate, los tokens tras la deduplicación, el bitext candidato, la precision humana y la distribución por dominio, en lugar de informar solo el tamaño total del crawl.
\end{warningbox}

\subsection{Los límites de zero-shot: una dirección no vista no es un idioma no visto}

NLLB no extrajo datos de entrenamiento para cada dirección $A\rightarrow B$; la explosión combinatoria lo haría costoso y no necesariamente acorde con las necesidades reales. Una dirección puede ser zero-shot, pero si el modelo ya vio los idiomas $A$ y $B$ en otras direcciones, todavía puede transferir conocimiento mediante las representaciones compartidas del encoder/decoder.

\begin{importantbox}{Hay que distinguir dos tipos de «no visto»}
\begin{itemize}
  \item \textbf{Unseen direction}: el entrenamiento no incluyó directamente $A\rightarrow B$, pero sí tareas como $A\leftrightarrow E$ y $B\leftrightarrow E$.
  \item \textbf{Unseen language}: el texto, los tokens o la supervisión del idioma $A$ casi nunca aparecieron.
\end{itemize}
En clase, zero-shot se refiere sobre todo al primer caso. Presentarlo como capacidad de traducir un idioma completamente nuevo exagera el alcance de la generalización.
\end{importantbox}

\teachervoice{Fan explica que algunos zero-shot pairs funcionan muy bien porque el modelo ya vio el idioma de entrada y el de salida, aunque no esa combinación dirigida; en conjunto, estas direcciones siguen siendo más débiles, y no se trata de traducción sin ejemplos previos de un unseen language por completo.}

\subsection{Del NLLB de texto a SeamlessM4T}

Ante la pregunta sobre voz, Fan aclaró que NLLB en sí es principalmente traducción de texto; fue el trabajo posterior, SeamlessM4T, el que integró speech, transcription y text translation. Esta respuesta separa con claridad los límites del proyecto: los supuestos sobre datos del pipeline de texto no pueden extrapolarse directamente a idiomas que dependen sobre todo de la lengua oral.

\begin{knowledgebox}{Cadena de evidencia de la extensión de texto a voz}
Speech translation requiere audio, transcripciones, cobertura de hablantes y acentos, segmentación del habla y representaciones acústicas; el NLLB de texto trabaja sobre todo con pares de oraciones escritas. SeamlessM4T conecta ambos mediante representaciones multimodales compartidas, pero sigue necesitando nuevas evaluaciones de voz y nuevos análisis de seguridad; no basta con reutilizar las puntuaciones de texto de FLORES.
\end{knowledgebox}

\subsection{Una nueva lectura en la era de los foundation models}

Al final de la clase se discutió que, si el proyecto se reiniciara en una era de modelos base más potentes, podría usarse un foundation model grande con supervised fine-tuning, lo que reduciría la necesidad de diseñar sistemas especializados para cada tarea. Fan recordó que en los inicios del NLP la investigación en translation, summarization y QA estaba separada, mientras que ahora cada vez más tareas se apoyan en una base compartida.

\teachervoice{El juicio del docente es una tendencia, no un teorema: cuanto más potente sea el modelo base, menos fine-tuning de dominio podría requerirse, pero sigue siendo una condición previa que los idiomas de bajos recursos estén bien representados en el preentrenamiento. Una base de propósito general no corrige por sí sola el desequilibrio de idiomas en el corpus de entrenamiento.}

\begin{warningbox}{No hay que borrar la contribución de NLLB con el relato posterior de los LLM}
Un foundation model más potente sigue necesitando datos multilingües, benchmarks confiables, language identification, evaluación humana y safety probes. El valor de largo plazo de NLLB está justamente en haber publicado de forma conjunta esa infraestructura y los modelos, no solo en haber propuesto una arquitectura de 2022.
\end{warningbox}

\subsection{Resumen de la sección}

Los sistemas multilingües futuros deben abarcar una representación explícita del idioma, las comunidades de la cola y barreras de ingeniería más bajas. La sesión de Q\&A añade precisiones: los datos y el modelo son igual de decisivos, la web abierta no está equilibrada, zero-shot se refiere sobre todo a direcciones no vistas, speech requiere evidencia de una nueva modalidad y un foundation model tampoco sustituye la infraestructura multilingüe.
\section{Síntesis y ampliación}

La contribución central de NLLB puede condensarse en una sola frase: \textbf{la traducción multilingüe a gran escala es un trabajo de ingeniería de sistemas que diseñan en conjunto las necesidades centradas en las personas, la infraestructura de evaluación, la producción de datos, la capacidad del modelo, la planificación del entrenamiento y la validación de seguridad.} Si solo se recuerda MoE o el 44\% de BLEU, se pierde la metodología que el curso realmente quiere transmitir.

\subsection{Doce conclusiones aplicables}

\begingroup
\footnotesize
\setlength{\columnsep}{1.2em}
\begin{multicols}{2}
\begin{enumerate}[itemsep=0.08em,topsep=0.15em,parsep=0pt,leftmargin=*]
  \item Primero confirmar las necesidades de la comunidad; después decidir idiomas, sistemas de escritura, direcciones y escenarios.
  \item Definir high/low-resource como una condición de los datos, no como un orden por población o por valor.
  \item Antes de entrenar, construir un evaluation set comparable, traducido por personas y que abarque varias direcciones.
  \item Reservar el costo real de los estándares lingüísticos, los traductores, los reviewers y el post-editing.
  \item Usar una pequeña cantidad de seed confiable para calibrar LID, el encoder, los filter y la MT inicial.
  \item Diseñar el web mining como un pipeline iterativo y rastreable, no como una extracción única.
  \item Registrar para cada conjunto de datos su provenance, noise, scale y teacher dependency.
  \item El balance de la cola larga debe aumentar la diversidad efectiva; no basta con duplicar las muestras de bajos recursos.
  \item La ganancia de capacidad de MoE debe evaluarse junto con el sobreajuste en bajos recursos, el enrutamiento y la regularización.
  \item Reportar los resultados por direction, resource group y metric, no solo el promedio.
  \item La evaluación automatic, la human y la de toxicity responden preguntas distintas.
  \item Solo al abrir el benchmark, las herramientas de datos y los modelos puede la comunidad seguir corrigiendo las brechas de cobertura.
\end{enumerate}
\end{multicols}
\endgroup

\subsection{Panorama final del sistema}

Al unir toda la clase se obtiene un ciclo de ingeniería con retroalimentación. Este ciclo no es un proceso lineal en cascada, porque la evaluación humana, las fallas de seguridad y la retroalimentación de la comunidad obligan al equipo a regresar a las etapas de estándares, datos y modelo:
\[
\begin{aligned}
\text{Community needs}
&\rightarrow\text{Standards \& Evaluation}
\rightarrow\text{Seed Data}
\rightarrow\text{Mining / BT}\\
&\rightarrow\text{MoE Training}
\rightarrow\text{Human \& Safety Validation}
\rightarrow\text{Community feedback}.
\end{aligned}
\]
Cada flecha puede introducir un sesgo de selección y cada retroalimentación puede modificar los estándares de las etapas previas. Un sistema verdaderamente maduro no se completa con un solo entrenamiento: es capaz de explicar de dónde provienen los datos, por qué el modelo falla en ciertos idiomas, qué métricas respaldan las conclusiones y quién tiene la facultad de corregir los errores.

\begin{importantbox}{Criterio de juicio final}
«No Language Left Behind» no puede declararse terminado al alcanzar cierto número de idiomas. Un juicio más confiable consiste en preguntar: ¿los idiomas de la cola cuentan con datos y evaluaciones auditables?, ¿los hablantes nativos participan en las necesidades, los estándares y la validación?, ¿el sistema reporta la calidad por dirección y los errores de alto riesgo?, ¿los recursos abiertos permiten que la comunidad reproduzca, critique y mejore? La cifra de cobertura solo puede ser una de las piezas de esta evidencia.
\end{importantbox}

\subsection{Lecturas adicionales}

\begingroup
\scriptsize
\begin{itemize}[itemsep=0pt,topsep=0.15em,parsep=0pt]
  \item NLLB Team, \href{https://arxiv.org/abs/2207.04672}{No Language Left Behind: Scaling Human-Centered Machine Translation}: informe completo sobre datos, modelos, evaluación e impacto social.
  \item Meta AI, \href{https://research.facebook.com/publications/no-language-left-behind/}{NLLB publication page} y \href{https://github.com/facebookresearch/fairseq/tree/nllb}{fairseq NLLB release}: punto de entrada oficial al proyecto y al código de los modelos.
  \item \href{https://github.com/facebookresearch/flores}{FLORES}: datos de evaluación multilingüe; \href{https://github.com/facebookresearch/stopes}{Stopes}: herramienta escalable de minería de datos multilingües.
  \item Schwenk et al., \href{https://arxiv.org/abs/1907.05791}{WikiMatrix} y \href{https://github.com/facebookresearch/LASER}{LASER}: fundamentos de la minería de corpus mediante vectores de oraciones.
\end{itemize}
\endgroup

\end{document}
